\section{The full distributional collision correction}
\label{cont:sec}

The incoming coincident-point report isolates the off-point collision
singularity and its finite constant, but leaves the distribution
supported at the collision undetermined \cite{IncomingCoincident}.
We now specify the extension convention and calculate that distribution
for every pair of Stieltjes indices and every pair of argument
derivatives.  This is a comparison of two extensions of the same
off-point function.  It does not identify a pointwise product of two
singular distributions at the same argument.

\subsection{The two extensions being compared}

Work on $\mathbb T=\mathbb R/\mathbb Z$.  Let $G_m$ be the unit-coordinate
Hadamard extension of $\gamma_m(x)$ from $0<x<1$, and let
$F_{m,p}$ be the corresponding extension of
$\gamma_m^{(p)}(x)$.  Thus, for a smooth periodic test function,
\begin{equation}\label{cont:finite-part}
 \langle F_{m,p},\phi\rangle
 =\CT_{\epsilon\downarrow0}
      \int_\epsilon^1\gamma_m^{(p)}(x)\phi(x)\,dx,
 \qquad F_{m,0}=G_m.
\end{equation}
Write $D=d/dx$, $\Delta=\delta_0-1$, and
$\check U(x)=U(-x)$.  The exact correlation is the well-defined
periodic convolution
\begin{equation}\label{cont:correlation}
 \mathcal R_{mn}^{p,q}=\check F_{m,p}*F_{n,q}.
\end{equation}
For ordinary smooth functions this is
$\int_0^1 f(x)g(x+a)\,dx$ as a function of $a$.

For $0<a<1$, let $J_{mn}^{p,q}(a)$ be the unit-coordinate finite part
of the separated product
$\int_0^1\gamma_m^{(p)}(x)\gamma_n^{(q)}(\{x+a\})\,dx$.
It is the off-point restriction of \eqref{cont:correlation}, as
established in the manuscript's translated contact law
\cite{Repo}.  Define its canonical extension in the shift coordinate by
\begin{equation}\label{cont:shift-extension}
 \left\langle\mathfrak F J_{mn}^{p,q},\phi\right\rangle
 =\CT_{\epsilon,\eta\downarrow0}
      \int_\epsilon^{1-\eta}J_{mn}^{p,q}(a)\phi(a)\,da.
\end{equation}
The two endpoint constants are taken independently in $a$ and $1-a$;
there is no added delta term.  The local power--log expansions proved
in the incoming collision report guarantee existence.

\subsection{Finite generating polynomials}

This section uses a raising parameter to simplify the Gamma quotient:
\[
 \mathcal G(u)=\sum_{m\ge0}(-1)^mG_m\frac{u^m}{m!},
 \qquad \mathcal F_p(u)=\sum_{m\ge0}(-1)^mF_{m,p}\frac{u^m}{m!}.
\]
For $r\ge0$ set
\begin{equation}\label{cont:P-H}
 P_r(z)=\prod_{j=1}^r(1+z/j),\qquad
 H_r(z)=\frac{P_r(z)-1}{z},\qquad P_0=1,\ H_0=0.
\end{equation}
$H_r$ is a polynomial, not the scalar harmonic number; its constant
coefficient is $\sum_{j=1}^r1/j$.  All its coefficients are elementary
symmetric polynomials in $1,1/2,\ldots,1/r$.

Put $w=u+v$ and define
\begin{equation}\label{cont:A-B}
 A(u,v)=\frac{\Gamma(1-u)\Gamma(1+w)}{\Gamma(1+v)},
 \qquad T_0(u,v)=\frac{A(u,v)-1}{uw},
 \qquad B(u,v)=T_0(u,v)+T_0(v,u).
\end{equation}
The quotients are removable because $A(0,v)=A(u,-u)=1$.
The convergent local series
\begin{equation}\label{cont:log-A}
 \log A(u,v)=\sum_{k\ge2}\frac{\zeta(k)}{k}
       \left[u^k+(-1)^k\big((u+v)^k-v^k\big)\right]
\end{equation}
makes every coefficient an exact finite polynomial in positive
integer zeta values.  The undifferentiated off-point use of this
Gamma quotient is inherited from the manuscript \cite{Repo}.

There is also a compact symmetric form:
\begin{equation}\label{cont:B-symmetric}
 B(u,v)=\frac1{uv}\left\{
  \frac{\Gamma(1-u)\Gamma(1-v)}{\Gamma(1-u-v)}
  \frac{\cos(\pi(u-v)/2)}{\cos(\pi(u+v)/2)}-1\right\}.
\end{equation}
Indeed,
$uvB=(vA(u,v)+uA(v,u))/(u+v)-1$.
The Gamma reflection formula turns the fraction before the subtraction
into the Gamma quotient in \eqref{cont:B-symmetric} times
$\{\sin(\pi u)+\sin(\pi v)\}/\sin(\pi(u+v))$.
The sine addition formula gives the displayed cosine ratio. All
apparent quotients are interpreted by their removable holomorphic
extensions near $(0,0)$. This second form provides an independent
coefficient check on the generator.

\begin{theorem}[All-index collision contact distribution]
\label{cont:main}
Let $m,n,p,q\ge0$ and $r=p+q$.  Then
\begin{equation}\label{cont:main-identity}
 \boxed{\mathcal R_{mn}^{p,q}
      =\mathfrak F J_{mn}^{p,q}+e_{mn}^{p,q}\delta_0^{(r)}}.
\end{equation}
There are no additional lower derivatives of $\delta_0$.
The coefficient is explicitly
\begin{equation}\label{cont:coefficient}
 e_{mn}^{p,q}=(-1)^{m+n}m!n![u^mv^n]E_{p,q}(u,v),
\end{equation}
where the following holomorphic germ requires only polynomial
arithmetic and \eqref{cont:log-A}:
\begin{align}
 E_{p,q}(u,v)=(-1)^p\bigg\{
 &P_r(w)B(u,v)
   +\frac{H_r(w)-H_r(v)}{u}
   +\frac{H_r(w)-H_r(u)}{v}\notag\\
 &-H_p(u)H_r(v)-H_q(v)H_r(u)+H_p(u)H_q(v)\bigg\}.
 \label{cont:E}
\end{align}
Both divided differences in this formula are polynomials.
\end{theorem}

\begin{proof}
Let $Z_s$ be the meromorphic periodic Hurwitz distribution.  Its
Fourier coefficients are the Hurwitz coefficients used in
Section~\ref{cub:sec}, with zero mean.  Local Taylor subtraction at
the origin gives
\begin{equation}\label{cont:Z-expansion}
 Z_{1+u}=-\frac\Delta u+\mathcal G(u).
\end{equation}
For $p>0$, the same subtraction at the higher Hurwitz pole, together
with $D^pZ_s=(-1)^p(s)_pZ_{s+p}$, gives
\begin{equation}\label{cont:derivative-anomaly}
 \mathcal F_p(u)=D^p\mathcal G(u)+H_p(u)\delta_0^{(p)}.
\end{equation}
This generating version of the manuscript's differentiation anomaly
also holds for $p=0$.  For clarity, the polar term of
$Z_{1+p+u}$ is $(-1)^{p+1}\delta_0^{(p)}/(p!u)$.
Multiplication by $(-1)^p(1+u)_p$ produces
$-P_p(u)\delta_0^{(p)}/u$; subtracting the pole of the left side
leaves exactly \eqref{cont:derivative-anomaly}.

The beta form of the classical correlation, continued as a periodic
distribution, is
\[
 \check Z_{1+u}*Z_{1+v}
 =-\frac{A(u,v)}u Z_{1+w}
   -\frac{A(v,u)}v\check Z_{1+w}.
\]
It follows either from the Fourier coefficients or first in a
convergent half-plane and then by meromorphic continuation.
Substitute \eqref{cont:Z-expansion} and use
$\Delta*\Delta=\Delta$ and $\Delta*G_m=G_m$.
The result is the holomorphic distribution identity
\begin{align}
 \check{\mathcal G}(u)*\mathcal G(v)
 ={}&B(u,v)\Delta
 -\frac{\mathcal G(w)-\mathcal G(v)}u
 -\frac{\check{\mathcal G}(w)-\check{\mathcal G}(u)}v
       \notag\\
 &-w\bigl[T_0(u,v)\mathcal G(w)
              +T_0(v,u)\check{\mathcal G}(w)\bigr].
 \label{cont:base-generator}
\end{align}
Its off-point constant is $-B$.  The canonical shift extension
replaces each ordinary Stieltjes function by its corresponding $G_m$
and retains that constant.  Consequently the difference for $r=0$
is $B\delta_0$, as claimed; it is not $B\Delta$.

For $r>0$, converting $D^r\mathcal G(z)$ to the canonical extension
of its ordinary derivative adds the difference
$-H_r(z)\delta_0^{(r)}$.  A reflected term has the same difference:
reflection and the chain rule each contribute $(-1)^r$.
Applying this to \eqref{cont:base-generator}, the contact coefficient
of its $r$th derivative is
\[
 P_r(w)B+
 \frac{H_r(w)-H_r(v)}u+
 \frac{H_r(w)-H_r(u)}v.
\]
Finally expand the correlation of
$D^p\mathcal G(u)+H_p(u)\delta_0^{(p)}$ and
$D^q\mathcal G(v)+H_q(v)\delta_0^{(q)}$.
The two cross terms and the delta--delta term supply, respectively,
$-H_p(u)H_r(v)$, $-H_q(v)H_r(u)$, and $H_p(u)H_q(v)$.
The orientation of the first factor contributes $(-1)^p$.
This proves \eqref{cont:E}.  Every contact term in this calculation
has derivative order $r$, proving the absence of lower orders.
Coefficient extraction gives \eqref{cont:main-identity}.
\end{proof}

\subsection{Explicit contact coefficients}

For undifferentiated factors, $E_{0,0}=B$ and the first values are
\begin{center}
\begin{tabular}{c c c}
\toprule
$(m,n)$ & $e_{mn}^{0,0}$ & contact term\\
\midrule
$(0,0)$ & $2\zeta(2)$ & $2\zeta(2)\delta_0$\\
$(1,0)$ or $(0,1)$ & $\zeta(3)$ & $\zeta(3)\delta_0$\\
$(2,0)$ or $(0,2)$ & $11\zeta(4)/2$ & $11\zeta(4)\delta_0/2$\\
$(1,1)$ & $7\zeta(4)/2$ & $7\zeta(4)\delta_0/2$\\
\bottomrule
\end{tabular}
\end{center}
The last two simplifications use $\zeta(2)^2=5\zeta(4)/2$.
For Stieltjes indices $m=n=0$ and positive argument derivatives,
\begin{align}
 e_{00}^{1,0}&=1-2\zeta(2),\notag\\
 e_{00}^{1,1}&=1-2\zeta(2),\notag\\
 e_{00}^{2,0}&=2\zeta(2)-\frac54.
 \label{cont:low-examples}
\end{align}
The associated distributions are $\delta_0'$, $\delta_0''$, and
$\delta_0''$, respectively.  Swapping the two factors reflects the
correlation; accordingly
\begin{equation}\label{cont:reflection}
 e_{nm}^{q,p}=(-1)^r e_{mn}^{p,q}.
\end{equation}
This identity is also immediate from \eqref{cont:E}.

For instance, the known off-point identity is
$J_{00}^{0,0}(a)=\gamma_1(a)+\gamma_1(1-a)-2\zeta(2)$.
Its exact periodic distributional completion is
\begin{equation}\label{cont:example00}
 \check G_0*G_0=G_1+\check G_1+2\zeta(2)(\delta_0-1).
\end{equation}
An off-point numerical computation cannot see the delta term.
In contrast, a nonzero Fourier coefficient of \eqref{cont:example00}
contains the necessary constant $2\zeta(2)$; omitting it gives the
wrong correlation even though the two functions agree for $0<a<1$.

\subsection{Normalized primitives of the contact discrepancy}

The contact formula also fixes constants in antiderivatives.
Let $\mathcal I$ be the zero-mean inverse of $D$ on periodic
distributions of mean zero.  For the periodic Bernoulli polynomial
$B_j(\{a\})$,
\begin{equation}\label{cont:Bernoulli-inverse}
 \mathcal I^j\Delta=-\frac{B_j(\{a\})}{j!},\qquad j\ge1.
\end{equation}
This follows from $DB_1(\{a\})=1-\delta_0$ and
$B_j'=jB_{j-1}$ for $j>1$, with zero means.
For $r>0$, applying $\mathcal I^{r+j}$ to the discrepancy in
\eqref{cont:main-identity} yields
\begin{equation}\label{cont:contact-primitive}
 \mathcal I^{r+j}\bigl(\mathcal R_{mn}^{p,q}
             -\mathfrak FJ_{mn}^{p,q}\bigr)
       =-e_{mn}^{p,q}\frac{B_j(\{a\})}{j!},\qquad j\ge1.
\end{equation}
For $r=0$, first subtract the mean of the discrepancy, replacing
$e\delta_0$ by $e\Delta$, and the same formula applies.
These are fixed normalized primitives, with no unspecified additive
polynomial.  They explain exactly how a point-supported correction
can become an ordinary Bernoulli term after repeated integration.
